\section{Second stage: Outcome model}
\label{sec:outcome-model}

In this section, we describe the second stage of the proposed two-stage estimation procedures. 
The  goal of the second stage is to estimate the causal mean response surface 
$\mathbb{E}[Y_i(\underline{a})]$ through a Marginal Structural Model (MSM). 
In this thesis, we focus on the truncated treatment histories similar to \cite{blackwellEffectPoliticalAdvertising2024}, 
where we model the effects of manipulating the last $k$ periods of treatment 
while treating earlier history as baseline confounding. 

\subsection{Common Framework and IPTW Identification}
Both the FE-MSM and SeqGPLVM methods use Inverse Probability of Treatment Weighting (IPTW) 
to construct a pseudo-population where the treatment history is independent of the measured and 
unmeasured time-invariant confounders. Under the assumption of sequential ignorability 
(either unit-specific or conditioned on a latent substitute), the truncated potential outcome 
$Y_i(\underline{a}_{T-k})$ is identified using the following weighted estimand:
$$\mathbb{E}[Y_i(\underline{a}_{T-k})] = \mathbb{E}\left[\frac{\mathbf{1}(\underline{A}_{i,T-k} = \underline{a}_{T-k})Y_i}{\prod_{t=T-k}^{T} \pi_{it}^{a_t}(1 - \pi_{it})^{1-a_t}}\right]$$
Weight Construction: We utilize **stabilized weights** to mitigate the high variance often associated with standard IPTW estimators. The denominator consists of the product of predicted propensities $\hat{\pi}_{it}$ from the first-stage treatment model.
Positivity and Trimming: To prevent numerically unstable or infinite weights, we apply **propensity truncation**. For both models, propensities are clipped at $\varepsilon = 10^{-6}$ to ensure they remain bounded away from 0 and 1. For units with no treatment variation (always-treated or never-treated), we follow Blackwell and Yamauchi’s approach of setting denominator propensities to $1-\varepsilon$ or $\varepsilon$ to avoid degeneracy.
Linear MSM Specification: We fit a linear, additive MSM to summarize the causal response surface:
    $$f(\underline{a}_{T-k}; \gamma) = \alpha' + \tau'_F a_T + \tau'_C \sum_{t=T-k}^{T-1} a_t$$
    where $\tau'_F$ represents the effect of the final treatment and $\tau'_C$ represents the cumulative effect of recent history.



\subsection{FE-MSM}
The Fixed-Effects MSM (FE-MSM) addresses time-invariant unmeasured confounding by incorporating unit-specific fixed effects directly into the propensity score model.

Unit-Specific Sequential Ignorability: This method relies on the assumption that treatment is randomized conditional on the observed history and a time-constant unit-specific effect $\alpha_i$.
Estimation Procedure: We obtain estimates for the propensity parameters $\beta$ and the unit effects $\alpha_i$ via **Conditional Maximum Likelihood**. This approach is justified under **large-$N$, large-$T$ asymptotics, where the bias from noisy fixed-effect estimation (the incidental parameters problem) fades over time.
Weighted Least Squares: Once the fixed-effect weights are constructed, the MSM parameters $\gamma$ are estimated by solving a weighted estimating equation, which reduces to **Weighted Least Squares (WLS)** for our linear specification. We use a **robust (sandwich) variance estimator** to obtain valid Wald confidence intervals.

\subsection{SeqGPLVM}

In contrast to the point-estimate approach of FE-MSM, the SeqGPLVM captures unmeasured confounding by inferring a posterior distribution $q(Z)$ over a latent substitute $Z_i$.

Handling Circular Dependency: To avoid the "circular dependency" created by using $Z_i$ as a covariate in a separate propensity model, we use the estimated propensities from the SeqGPLVM **directly** to build the IPTW weights.
Multiple Imputation (MI)-Style Procedure: To propagate the latent uncertainty from the GP-based treatment model, we draw $S$ samples $Z^{(s)} \sim q(Z)$. For each draw, we:
    1.  Compute draw-specific propensities and weights $W^{(s)}$.
    2.  Fit the weighted MSM to obtain a draw-specific estimate $\hat{\gamma}^{(s)}$ and robust covariance $V^{(s)}$.
Rubin’s Rules for Pooling: The final point estimate is the average across draws, $\bar{\gamma} = S^{-1}\sum \hat{\gamma}^{(s)}$. The total variance is calculated using **Rubin’s within/between decomposition**, $T = W + (1 + 1/S)B$, which explicitly accounts for the variability induced by the uncertainty in the inferred latent confounder.



Although the seqgplvm can be used in the settings in which the outcome is observed at every time point, 
since the asymptotic properties and consistency of their IPTW-FE estimator 
relies on having less number of N than T (assumption )
Based on the consistency assumption \ref{ass:C-consistency}, we can define 
truncated potential outcomes 